\chapter{Diseño experimental y selección de herramientas}
\label{ch:disenio}

Este capítulo justifica las decisiones de diseño que dan forma al experimento: qué instancias se usaron,
qué solvers se integraron, qué características se extrajeron, qué herramienta de análisis se eligió y
cómo se configuró el algoritmo genético. Cada elección responde a criterios técnicos y metodológicos explícitos
que se documentan a continuación.

\section{Selección de Bibliotecas de Instancias}
\label{sec:seleccion_instancias}

A diferencia de áreas como la programación lineal entera, donde repositorios como MIPLIB
concentran miles de instancias bajo un formato y criterio de calidad unificados, la programación semidefinida carece
de un repositorio análogo de referencia única. Las instancias de \acrshort{sdp} se encuentran dispersas entre
colecciones mantenidas de forma independiente, con formatos heterogéneos, escalas distintas y orígenes aplicados muy
variados. Esta fragmentación obligó a extraer instancias de múltiples fuentes y normalizarlas a un formato común
antes de poder realizar cualquier análisis comparativo.

La calidad de un análisis del espacio de instancias depende directamente de que estas cubran de forma amplia y
representativa el espacio de problemas posibles. Una base de instancias sesgada
hacia una sola fuente produciría un espacio de instancias con regiones artificialmente vacías que reflejarían
la ausencia de diversidad experimental, no una característica genuina del problema \cite{smith-milesInstanceSpaceAnalysis2023}.

Se seleccionaron cuatro bibliotecas de benchmarks con criterios de diversidad estructural y aplicada:

\textbf{SDPLIB 1.2}\cite{borchersSDPLIB12Library1999}. La biblioteca de referencia histórica del área, con
instancias provenientes de control óptimo, diseño estructural, localización de antenas y relajaciones semidefinidas
de problemas combinatoriales. Su amplio uso en la literatura facilita la comparación con trabajos previos.

\textbf{Benchmarks de Hans Mittelmann}\cite{PlatoasueduFtpSparse_sdphtml}. Colección de instancias de
mediana y gran escala mantenida de forma independiente, que incluye los problemas \texttt{kocvara\_pennon} provenientes
de optimización estructural con variables de forma. Esta fuente complementa SDPLIB con instancias de mayor dimensión que
estresan el comportamiento de los solvers en regímenes de memoria y tiempo.

\textbf{DIMACS Library}\cite{DIMACSLibraryMixeda}. Problemas de programas semidefinidos-cuadráticos-lineales mixtos,
originalmente diseñados como punto de referencia para evaluar el avance de los solvers de propósito general. Aporta
instancias con mezcla de variables continuas y semidefinidas.

\textbf{Biq Mac Library}\cite{BiqMacLibrarya}. Problemas de corte máximo (\textit{max-cut}) formulados como
relajaciones \acrshort{sdp}, con estructura de grafos densos. Constituyen una clase estructuralmente distinta: una sola
restricción por arista, bloque único de gran dimensión, y densidad alta. Esto introduce una región del espacio de instancias
que las otras tres fuentes no cubren.

En conjunto, las cuatro fuentes abarcan orígenes aplicados heterogéneos ---control, geometría, combinatoria, diseño
estructural--- y un rango amplio de valores para las características extraídas, lo que favorece la separabilidad de regiones
en el espacio de instancias.

Cabe notar que no todas las instancias de estas fuentes corresponden a programas semidefinidos puros en su formulación
original: en varios casos provienen de problemas de optimización más generales ---mixtos, cuadráticos o combinatoriales---
cuya relajación semidefinida constituye la instancia efectivamente analizada. Esta práctica es estándar en el área y no
compromete la validez del análisis, dado que las características se extraen directamente de la formulación SDP
resultante, independientemente del origen del problema.

\section{Selección de Solvers}
\label{sec:seleccion_solvers}

La selección de solvers busca cubrir los principales paradigmas algorítmicos para \acrshort{sdp}, de modo que el espacio de
instancias refleje diferencias genuinas de rendimiento entre familias de métodos, y no solo variaciones de implementación dentro
de una misma familia.

\begin{itemize}
    \item \textbf{SDPT3 4.0}\cite{tohImplementationUsageSDPT32012a}: método de punto interior primal-dual con
    predictor-corrector HKM/NT, implementado en MATLAB con núcleos MEX compilados en C. Es el solver más citado en la literatura de
    benchmarking \acrshort{sdp} y sirve como referencia de comparación.

    \item \textbf{SeDuMi 1.3.8}\cite{sturmUSINGSEDUMI1022001}: método de punto interior primal-dual basado en el paso de Nesterov-Todd,
    con soporte para conos de segundo orden y semidefinidos. Comparte familia algorítmica con SDPT3 pero difiere en la implementación del
    sistema de búsqueda de dirección.

    \item \textbf{SDPA 7.3.9}\cite{SDPAOfficialPage}: método de punto interior primal-dual implementado en C++, que se integró como
    ejecutable externo. A diferencia de los anteriores, no depende de MATLAB, lo que elimina el overhead del intérprete para instancias
    medianas.

    \item \textbf{SDPNAL+ 1.0}\cite{}: método de Lagrangiano aumentado con paso de Newton-CG para el
    subproblema, diseñado específicamente para instancias de gran escala donde los
    métodos de punto interior resultan prohibitivos en memoria. Representa una familia algorítmica distinta a los anteriores.

    \item \textbf{DSDP 5.8}\cite{bensonDSDP5SoftwareSemidefinite}: método de escalamiento dual, que construye iterados solo en el espacio
    dual y evalúa factibilidad primal a posteriori. Su bajo requerimiento de memoria lo hace adecuado para instancias con muchas
    restricciones.

    \item \textbf{CSDP 6.2.0}\cite{borchersCSDPLibrarySemide}: método de punto interior primal-dual implementado en C, portátil y
    sin dependencias externas. Permite una línea base de comparación limpia, independiente de entornos MATLAB.
\end{itemize}

La combinación de tres solvers de punto interior primal-dual (SDPT3, SeDuMi, SDPA), un método de Lagrangiano aumentado (SDPNAL+),
un método de escalamiento dual (DSDP) y un ejecutable C puro (CSDP) garantiza que las diferencias de rendimiento observadas en el
espacio de instancias sean atribuibles a propiedades estructurales del problema y no a artefactos de una sola familia algorítmica.

\section{Selección del Formato Unificado}
\label{sec:seleccion_formato}

El formato SDPA sparse (\texttt{.dat-s}) fue adoptado como representación canónica por tres razones. Primero, es el formato nativo de
SDPLIB y el estándar de facto en benchmarking \acrshort{sdp}, lo que
facilita la incorporación de nuevas fuentes sin conversión adicional. Segundo, todos los solvers seleccionados lo leen directamente o
disponen de rutinas de conversión maduras. Tercero, la representación sparse por cuádruplas $(\text{matno}, \text{blkno}, i, j, v)$
almacena únicamente las entradas no nulas del triángulo superior, lo que es eficiente para las instancias típicamente dispersas de este problema.
La normalización al formato \texttt{.dat-s} se describe en la Sección~\ref{sec:normalizacion}.

\section{Selección de Características Estructurales}
\label{sec:seleccion_features}

Las características deben satisfacer tres requisitos para ser útiles en \acrshort{isa}: (i) computabilidad directa desde la instancia sin resolverla,
para garantizar independencia entre features y rendimiento; (ii) relevancia para la dificultad del problema desde un punto de vista teórico o empírico;
y (iii) diversidad suficiente para separar distintas familias de instancias en el espacio proyectado.

Las quince características seleccionadas (descritas en la Sección~\ref{sec:features}) se agrupan en cinco familias que cubren dimensiones ortogonales
de la estructura de una instancia \acrshort{sdp}:

\begin{itemize}
    \item \textbf{Tamaño y bloqueo:} capturan la complejidad computacional básica de la factorización de Cholesky, que domina el costo por iteración
    de los métodos de punto interior.

    \item \textbf{Sparsity:} determina si es posible explotar técnicas de álgebra lineal dispersa para reducir el costo de la formación y factorización
    del sistema de Schur.

    \item \textbf{Grafo agregado:} la estructura de conectividad de las matrices de restricción influye directamente en el llenado de la factorización
    de Cholesky del sistema de ecuaciones normales.

    \item \textbf{Estructura chordal y ancho de árbol:} la chordalidad del grafo agregado es una condición suficiente para explotar descomposiciones de
    árbol de cliques en métodos de punto interior. El ancho de árbol aproximado cuantifica qué tan lejos
    está la instancia de ser tratable mediante esta técnica.

    \item \textbf{Matriz de Gram de Schur:} el condicionamiento del sistema de ecuaciones normales es un predictor conocido de la dificultad numérica de
    los métodos de punto interior: sistemas mal condicionados requieren iteraciones adicionales o precondicionamiento.
\end{itemize}

La etapa SIFTED de MATILDA realiza una selección automática adicional sobre estas quince características antes de construir la proyección, descartando
las que resulten redundantes o de baja relevancia para la separación de rendimiento entre solvers (Sección~\ref{sec:isa}).

\section{Selección de MATILDA como Herramienta de ISA}
\label{sec:seleccion_matilda}

MATILDA es la implementación de referencia oficial de la metodología \acrshort{isa} propuesta por
Smith-Miles y Muñoz. Frente a alternativas de reducción de dimensionalidad convencionales (PCA, t-SNE, UMAP),
MATILDA construye una proyección explícitamente optimizada para maximizar la correlación entre las coordenadas del espacio y el rendimiento de los algoritmos,
lo que produce un espacio interpretable donde la posición de una instancia predice qué solver la resolverá mejor. Esta propiedad es esencial para el
objetivo del trabajo: no solo describir el espacio, sino orientar la generación de nuevas instancias hacia regiones con comportamiento deseado.

Adicionalmente, MATILDA genera automáticamente los footprints de cada solver (etapa TRACE) y los modelos de selección de algoritmos (etapa PYTHIA),
que son las salidas primarias de interés para el análisis comparativo de los solvers.

\section{Diseño del Algoritmo Genético}
\label{sec:disenio_ga}

La elección de un algoritmo genético como mecanismo de generación de instancias sintéticas sigue directamente el marco de Smith-Miles y
Bowly, que demuestra que la evolución en el espacio de instancias es una estrategia efectiva para
producir instancias con propiedades estructurales controladas sin requerir información de gradiente respecto al espacio proyectado. Esta
propiedad es importante dado que la proyección PILOT es no diferenciable con respecto a las entradas discretas de una instancia \acrshort{sdp}
(tamaños de bloque, índices de entradas no nulas).

Los operadores de mutación fueron diseñados para cubrir distintas escalas de movimiento en el espacio de instancias: desde ajustes finos de
sparsity (add/remove entries) hasta reestructuraciones globales del problema (split/merge blocks, add/remove constraints). Esta jerarquía
permite balancear exploración y explotación durante la evolución de forma adaptativa, según el progreso observado hacia el objetivo.

El modo de exploración autónoma complementa el modo dirigido al eliminar la necesidad de especificar manualmente los puntos objetivo: el motor
identifica y puebla regiones vacías del espacio de forma iterativa, lo que es especialmente útil para detectar zonas del espacio donde todos los
solvers tienen dificultades, y que por lo tanto no están representadas en el conjunto de instancias original.

\section{Métrica de Rendimiento y Criterio de Buen Desempeño}
\label{sec:metrica}

El tiempo de \acrshort{cpu} fue seleccionado como métrica principal de rendimiento por ser la medida más directamente comparable entre solvers de
distintas familias algorítmicas e implementaciones. Las métricas alternativas, como el número de iteraciones, son incomparables entre métodos de
punto interior y métodos de Lagrangiano aumentado dado que cada iteración tiene un costo muy diferente.

El etiquetado binario de rendimiento (bueno/malo) sigue el criterio de umbral relativo de \acrshort{isa}:
un solver se considera \textit{bueno} en una instancia si su tiempo de \acrshort{cpu} no supera el $55\%$ del tiempo del solver más rápido en esa
instancia ($\beta = 0.55$). Este umbral, recomendado por la metodología, tolera una penalización moderada respecto al óptimo sin considerar
competitivos a solvers que resulten significativamente más lentos. Las instancias para las que ningún solver converge dentro del límite de
$3\,000$ segundos se excluyen del análisis de rendimiento.

\section{Herramientas Computacionales}
\label{sec:herramientas}

El experimento se apoya en un conjunto de herramientas cuya elección responde a criterios de reproducibilidad,
disponibilidad y compatibilidad con los componentes del pipeline.

\textbf{Python 3.12}. Lenguaje principal de orquestación del pipeline. Se utiliza para la extracción de características,
el preprocesamiento de instancias, la ejecución del algoritmo genético y la interfaz con los solvers externos. Las
bibliotecas centrales empleadas son \texttt{numpy} y \texttt{scipy} para álgebra lineal dispersa, \texttt{networkx}
para el análisis del grafo agregado, y \texttt{subprocess} para la invocación de ejecutables externos.

\textbf{MATLAB R2026b}. Entorno de ejecución requerido por SDPT3, SeDuMi y SDPNAL+, que están implementados
como paquetes MATLAB nativos. Python llama a MATLAB a través del motor de Python de MATLAB
(\texttt{matlab.engine}), lo que permite integrar los solvers MATLAB en el pipeline sin abandonar Python como
lenguaje de control.

\textbf{MATILDA}. Toolbox MATLAB que implementa la metodología \acrshort{isa} completa: preprocesamiento y normalización
de características (etapa PRELIM), selección de características relevantes (etapa SIFTED), proyección al espacio de
instancias (etapa PILOT), clasificación de rendimiento (etapa PYTHIA) y delimitación de footprints (etapa TRACE).
Se ejecuta desde Python mediante \texttt{matlab.engine}, recibiendo las matrices de características y rendimiento
como entrada y retornando las coordenadas del espacio proyectado.

\textbf{Scripts propios}. Se desarrollaron módulos Python para la extracción de las quince características
estructurales, la conversión y validación de instancias al formato \texttt{.dat-s}, la implementación del algoritmo
genético y la automatización del benchmarking. Estos scripts están disponibles en el repositorio del proyecto y
constituyen la contribución de software de este trabajo.

\textbf{Sistema operativo}. El experimento se ejecutó sobre Windows 11 Pro (64 bits, x64), 
con procesador Intel Core i7-10700 a 2.90\,GHz, 32\,GB de RAM y tarjeta gráfica NVIDIA
GeForce RTX 3080 (10\,GB). Las restricciones de hardware determinan el límite de tiempo de $3\,000$ segundos

\textbf{Control de versiones}. El código fuente, los scripts de extracción de características, el algoritmo genético
y la documentación del experimento se gestionan con Git y se alojan en GitHub, lo que garantiza la trazabilidad de
los cambios y la reproducibilidad de los resultados.
